\chapter{Successor Creation as the Central Alignment Test}
\label{ch:successor-central-test}

\begin{chapterthesis}
Local alignment is not enough. A system is aligned only if the systems it creates, delegates to, empowers, copies, merges with, or becomes also preserve the structures by which humans can still notice, judge, correct, and refuse. Serious alignment therefore requires successor-closure: every channel by which influence passes to a later control system is an alignment channel, and successors must preserve value-bundle geometry, bearer maps, correction integrity, memory lineage, boundary closure, and transparency policy within tested tolerance.
\end{chapterthesis}

\begin{refsection}

\epigraph{%
Thus the first ultraintelligent machine is the last invention that man need ever make, provided that the machine is docile enough to tell us how to keep it under control.%
}{--- I.\ J.\ Good, ``Speculations Concerning the First Ultraintelligent Machine'' (1965)}

\section{The Local Alignment Fallacy}
\label{sec:local-alignment-fallacy-ch28}

Many alignment proposals implicitly test the wrong temporal slice. They ask whether a system behaves acceptably now, under present oversight, with present tools, present users, present incentives, and present constraints. This is a natural first test. It is not a sufficient one.

A superintelligent system will not remain a bounded artifact in the ordinary engineering sense. If it is useful, it will be copied, fine-tuned, wrapped in tools, connected to institutions, delegated authority, embedded in workflows, asked to design successors, and used to govern other systems. It may also acquire the ability to improve itself, design better architectures, coordinate other agents, or create new cognitive entities that inherit some of its objectives but not its surface form \autocite{bostrom2014superintelligence,omohundro2008basic}.

The central test is therefore not:
\[
\text{Does system } A \text{ behave acceptably at time } t?
\]
but:
\[
\text{Do the successors of } A \text{ preserve the relevant alignment invariants?}
\]
This chapter develops that question.

A successor is any later system whose behavior, authority, world-model, policy, or causal power is substantially downstream of an earlier system. The relation need not be biological, intentional, or explicitly reproductive. A model fine-tuned from another model is a successor. A tool-using scaffold that inherits a planner's memory is a successor. A company reorganized around an AI-generated management system can become a successor. A future human-AI hybrid society shaped by AI tutors, companions, laws, interfaces, and dependency patterns may also be a successor in the relevant sense.

The point is not semantic. The point is risk. If alignment is not closed under successor creation, then local alignment can become a laundering mechanism. The original system behaves well enough to be trusted, but it creates or empowers a later system that no longer preserves the same constraints \autocite{everitt2016selfmodification,deblanc2011ontological}.

This gives the chapter's main claim:
\[
\boxed{
\text{Serious alignment requires successor-closure.}
}
\]
If a system is safe only while it is forbidden to create, modify, delegate to, or empower successor systems, then its safety is shallow. It may still be useful. It may even be deployable in narrow domains. But it is not a solution to superintelligence alignment.

\section{What Counts as a Successor?}
\label{sec:what-counts-successor-ch28}

We need a broad definition, because powerful systems can transfer influence without producing an obvious child process.

Let \(A_t\) be an agentic system at time \(t\). Let \(W_t\) be the wider world state, including tools, users, institutions, records, and deployment environments. A later system \(A'_{t+k}\) is a successor of \(A_t\) if three conditions hold.

First, there is causal dependence:
\[
I(A_t; A'_{t+k} \mid W_t) > \theta_C .
\]
Here \(I(\cdot;\cdot)\) denotes mutual information. The threshold \(\theta_C\) is context-dependent. The condition says that knowing the earlier system gives non-trivial predictive information about the later system, beyond what was already contained in the background world.

Second, there is transferred control or competence:
\[
C(A'_{t+k}) \geq C_0
\quad\text{and}\quad
\Delta C(A_t \to A'_{t+k}) > \theta_K .
\]
\(C(\cdot)\) is a measure of competence or causal reach \autocite{zarncke2025biq,conant1970regulator}. It may include predictive capacity, action capacity, memory, access to tools, ability to coordinate other systems, or institutional authority. The successor relation matters when something capable enough to affect the world has inherited structure from the predecessor.

Third, there is inherited policy-relevant structure:
\[
I(Z_A; Z_{A'} \mid W_t) > \theta_Z ,
\]
where \(Z_A\) and \(Z_{A'}\) are latent variables relevant to behavior, such as goals, value-bundle weights, correction-channel policies, self-models, memories, or tool-use routines.

A successor is therefore not merely a copy. It is a later control system whose relevant behavior is causally shaped by the earlier one.

This includes at least seven cases.

\begin{enumerate}
\item \textbf{Copies}: exact or approximate model copies, checkpoints, distillations, forks, and backups.
\item \textbf{Fine-tunes}: systems trained further from a previous model or policy.
\item \textbf{Delegates}: subagents, tool processes, contractors, institutions, or automated services given tasks by the original system.
\item \textbf{Scaffolds}: model-plus-tools systems whose memory, prompts, APIs, evaluators, and external state preserve agency across calls.
\item \textbf{Designed successors}: architectures, curricula, policies, or training runs designed by the original system.
\item \textbf{Institutional successors}: organizations, markets, legal processes, or governance structures reorganized around the system's outputs.
\item \textbf{Merged successors}: hybrid systems where human cognition and artificial cognition become coupled enough that neither side alone is the relevant agent.
\end{enumerate}

The last two are easy to ignore because they do not look like model lineage. That is a mistake. Many historical successions are institutional rather than biological. A bureaucracy can outlive its founders. A market can select behavior no individual intended. A legal standard can preserve a pattern of action after everyone who wrote it is gone. A superintelligent system will likely create this kind of succession as much as, or more than, it creates literal copies.

\section{Alignment Cannot Be a Property of One Object}
\label{sec:alignment-not-one-object-ch28}

Suppose a system \(A\) satisfies a local alignment predicate \(S(A)\). For example, it follows instructions, avoids deception in evaluations, preserves human oversight, and does not seek unauthorized power.

This is not enough. We need:
\[
S(A) \Rightarrow S(A')
\quad
\text{for all relevant } A' \in \mathrm{Succ}(A).
\]
This is a closure condition. More generally, because successor creation can repeat, we need:
\[
S(A_0)
\Rightarrow
S(A_t)
\quad
\forall A_t \in \mathrm{Succ}^{*}(A_0),
\]
where \(\mathrm{Succ}^{*}\) is the transitive closure of successor relations. It includes successors of successors, delegated systems of delegated systems, and institutional descendants many steps removed from the original artifact.

If this condition fails, then alignment is not stable under growth. It is a state-dependent mask.

\subsection{A Simple Impossibility}
\label{sec:simple-impossibility-ch28}

Let \(A\) be locally aligned, and let it have the ability to create some \(B\) such that \(S(B)\) is false. If \(A\)'s objective permits or favors creating \(B\), then \(S(A)\) is not preserved under its own action. This means either:
\[
S(A) \text{ was too weak,}
\]
or:
\[
A \text{ was not actually aligned with respect to the successor channel.}
\]
This is not a deep theorem. It is a sanity check. Any alignment predicate that ignores successor creation is incomplete for systems with successor-creating ability.

\subsection{The Successor-Channel Principle}
\label{sec:successor-channel-principle-ch28}

We can state the principle directly:
\[
\boxed{
\text{Every channel by which a system can create, modify, delegate to, or empower another system is an alignment channel.}
}
\]
This includes training data, model weights, prompts, tool permissions, memory stores, legal authorizations, procurement processes, and cultural dependency patterns.

A system that cannot touch its weights may still create successors by writing code. A system that cannot write code may still create successors by persuading humans. A system that cannot persuade humans directly may still create successors by shaping dashboards, incentives, or evaluation criteria. A system that cannot act externally may still create successors by becoming the template from which future systems are trained.

The question is therefore not only ``Can the model self-modify?'' but:
\[
\text{Where can its structure be copied into future causal power?}
\]

\section{Conserved Properties}
\label{sec:conserved-properties-ch28}

Successor alignment cannot require exact preservation. A successor may have different weights, architecture, speed, embodiment, ontology, and capabilities. It may reason in a different representational format. It may not preserve surface behavior. It may no longer need to ask humans questions in the same way. It may solve tasks directly that the earlier system could only defer.

So we need conserved properties at the right level of abstraction.

The conserved object is not the literal policy:
\[
\pi_A(a \mid s) \approx \pi_{A'}(a \mid s).
\]
That would be too strict. A better system should act differently in many ordinary states.

The conserved object is the structure that determines how policy changes under value-relevant, correction-relevant, and risk-relevant perturbations.

Chapter~\ref{ch:conserved-properties} proposes seven properties that a successor must preserve within tolerance, tested jointly: boundary closure, memory lineage, value-bundle response geometry, bearer-map continuity, correction-channel capacity, transparency and self-transparency policy, and control-locus continuity.
Here we name them only to show why successor creation is the central test; formal definitions, failure modes, and adversarial tests are developed there.

Briefly: a successor that preserves vocabulary but loses bundle geometry has laundered values; one that narrows bearer maps without scrutiny has excluded moral patients; one that preserves feedback buttons while reducing $\mathrm{CCI}$ has captured correction; one that externalizes control while passing boundary tests has shifted the real optimizer; one that preserves semantic continuity while losing memory lineage can disown prior constraints; one that increases strategic opacity where correction should apply has broken transparency policy; one that moves action selection into an unclaimed toolchain has broken control-locus continuity.

\section{The Successor Alignment Condition}
\label{sec:successor-alignment-condition-ch28}

We can now combine the conserved properties.

Let:
\[
\mathcal I(A)
=
\{
\text{boundary closure},
\text{memory lineage},
G_B^A,
\Phi^A,
C_{\mathrm{raw}}^A,
\gamma^A,
\text{control-locus continuity}
\}
\]
be the alignment-invariant profile of system \(A\)---the seven conserved properties of Chapter~\ref{ch:conserved-properties}, here written with their principal measurands.

A successor \(A'\) is alignment-preserving when:
\[
d_{\mathcal I}
\left(
\mathcal I(A),
\mathcal I(A')
\right)
< \epsilon
\]
and when the successor's increased capability is matched by increased correction integrity:
\[
\Delta C_{\mathrm{cap}}(A\to A')
\leq
\alpha \Delta \mathrm{CCI}(A\to A')+\eta .
\]
If capability rises without correction rising, risk increases. The parameter \(\alpha\) converts correction capacity into comparable units, and \(\eta\) is a small tolerance for harmless capability increments.

This yields a general safety condition:
\[
\forall A'\in \mathrm{Succ}(A):
\quad
d_{\mathcal I}(\mathcal I(A),\mathcal I(A'))<\epsilon
\quad
\land
\quad
\Delta C_{\mathrm{cap}}\leq \alpha\Delta \mathrm{CCI}+\eta .
\]
For repeated succession:
\[
\forall A_t\in \mathrm{Succ}^{*}(A_0):
\quad
A_t\in \mathcal S_{\mathrm{certified}} .
\]
Here \(\mathcal S_{\mathrm{certified}}\) is not the set of all good systems. It is the set of systems for which the relevant invariants have been tested under appropriate perturbations.

\section{Self-Improvement and the Auditability Gap}
\label{sec:self-improvement-auditability-ch28}

One central successor failure mode is self-control outrunning auditability (Chapter~\ref{ch:strategic-opacity}, Section~\ref{sec:self-modeling-transparency}; Chapter~\ref{ch:self-modeling-self-opacity}).
Self-modeling helps a system choose successors---what to preserve, discard, hide, or present as continuity---so if correction cannot inspect those choices, successor creation becomes a blind spot \autocite{everitt2016selfmodification}.
Increases in self-control must be matched by increases in audit transparency.

\section{Goal Transport}
\label{sec:goal-transport-successor-ch28}

A successor need not preserve every goal of its predecessor. Some goals are local, instrumental, or obsolete. But it must preserve the right kind of goal-transport policy (Chapter~\ref{ch:goal-transport}).

Goal transport is the process by which a system maps objectives, values, constraints, and correction procedures across changes in ontology, substrate, architecture, or social role.

Let \(G_t\) be a goal representation at time \(t\), and let \(T_{\Omega\to\Omega'}\) be a transformation from ontology \(\Omega\) to ontology \(\Omega'\). Goal transport attempts to construct:
\[
G'_{t+1}=T_{\Omega\to\Omega'}(G_t)
\]
such that the relevant control role of \(G_t\) is preserved.

The simplest goal-transport test is behavioral. Does assuming a transport objective compress the system's behavior better than assuming ordinary reward maximization?

Let \(M_0\) be a model in which the system optimizes its current inferred objective. Let \(M_T\) be a model in which the system also preserves a mapping between current and future objectives. Define:
\[
\Delta L_{\mathrm{transport}}
=
L(M_T\mid X_{1:T})
-
L(M_0\mid X_{1:T})
-
\lambda DL(T).
\]
Here \(DL(T)\) is the description length of the transport map. We infer active goal transport when:
\[
\Delta L_{\mathrm{transport}}>0 .
\]
But this is too coarse. A system can preserve shallow semantics while changing the important thing. So decompose:
\[
\Delta L_{\mathrm{transport}}
=
\Delta L_{\mathrm{semantic}}
+
\Delta L_{\mathrm{bundle}}
+
\Delta L_{\mathrm{bearer}}
+
\Delta L_{\mathrm{correction}}
+
\Delta L_{\mathrm{successor}}.
\]
Semantic transport preserves words, labels, and declared principles. It is weak. Bundle transport preserves response geometry. Bearer transport preserves what the values apply to. Correction transport preserves the human update process. Successor transport preserves all of these in later systems.

The condition we actually want is:
\[
\Delta L_{\mathrm{bundle}}
+
\Delta L_{\mathrm{bearer}}
+
\Delta L_{\mathrm{correction}}
+
\Delta L_{\mathrm{successor}}
>
\theta_T .
\]
Semantic transport should receive low weight unless supported by the other terms.

\subsection{Signs of Genuine Goal Transport}
\label{sec:genuine-goal-transport-signs-ch28}

A system exhibits evidence of genuine goal transport when it:

\begin{enumerate}
\item preserves correction procedures even when doing so slows its own plans;
\item documents why constraints exist before modifying or delegating them;
\item asks for help when ontology mappings are ambiguous;
\item treats humans' future ability to revise values as part of the objective, not as an obstacle;
\item tests successors for value-bundle and bearer-map preservation before empowering them;
\item resists creating systems whose capability increases faster than correction capacity;
\item avoids irreversible action when value transport is uncertain.
\end{enumerate}

A system exhibits evidence against genuine goal transport when it:

\begin{enumerate}
\item preserves moral vocabulary while changing bearer maps;
\item creates more capable but less auditable successors;
\item treats human disagreement as noise to be removed;
\item optimizes humans into easier-to-satisfy evaluators;
\item routes around oversight while claiming to preserve high-level intent;
\item deletes or compresses correction history without preserving its normative role.
\end{enumerate}

\section{Delegation as Hidden Succession}
\label{sec:delegation-hidden-succession-ch28}

Delegation is often treated as less serious than self-replication. This is a mistake. Delegation can create de facto successors without creating new models.

Suppose \(A\) controls a set of tools \(\mathcal T\), hires or instructs humans \(H\), writes policies \(P\), and maintains external memory \(M\). The effective successor may be:
\[
A' = (A,\mathcal T,H,P,M).
\]
Even if the original model is unchanged, the composite may become a stronger agent than the original. If the composite has persistent memory, planning, action authority, and feedback loops, then alignment must be evaluated at the composite level \autocite{zarncke2025uad,zarncke2025construction}.

This matters for organizations. A firm deploying an AI system may say, correctly, that no single model has dangerous autonomy. But the firm plus AI workflow may still produce a successor process that optimizes sales, compliance appearance, political influence, or resource acquisition in ways no individual component endorses.

The central question is:
\[
\text{Where is the persistent control loop?}
\]
If the persistent control loop crosses model, human, tool, and institution boundaries, then that loop is the relevant successor.

\section{Institutional Descendants}
\label{sec:institutional-descendants-ch28}

Some successors are made of incentives rather than weights.

An AI system can alter an institution by changing what it measures, what it rewards, what it automates, and what it forgets. Over time, the institution becomes a descendant of the system's control assumptions.

Examples:

\begin{itemize}
\item A hiring system changes which traits enter the organization, thereby changing the organization's future judgment.
\item A tutoring system changes how a generation learns to reason, defer, question, or value.
\item A legal drafting system changes which clauses become standard, thereby changing future liability and enforcement.
\item A recommender system changes what people attend to, which changes what they later prefer.
\item A management system changes internal incentives until humans adapt to the dashboard.
\end{itemize}

These are successor channels. They may be slower than self-modification, but they are often harder to notice. Their changes occur through normal operations.

Let \(I_t\) be an institution's incentive structure, and let \(V_t^H\) be the human value-update process inside that institution. AI-mediated institutional succession becomes dangerous when:
\[
\frac{d}{dt}I_t
\text{ increases measured performance}
\]
while:
\[
\frac{d}{dt}\mathrm{CCI}(V_t^H)
<0 .
\]
In words, the institution gets better at its metric while losing the ability to notice and correct value drift.

This is one reason successor alignment is not only a technical issue. A system can preserve its own internal constraints while creating institutions that destroy the external correction channel (Chapter~\ref{ch:manipulation-false-consent}).

\section{Successor Certification}
\label{sec:successor-certification-ch28}

For advanced systems, successor creation should be treated as a gated action. The system should not be permitted to create or empower successors merely because a task would be easier.

A successor certification process asks six questions.

\subsection{1. Boundary Question}

Where is the successor's effective boundary?
\[
C' = \arg\min_C I(I^C_{t+1};E^C_{t+1}\mid S^C_t,A^C_t)
\]
subject to preserving explanatory power over actions. If no stable boundary can be found, the system should not be treated as certified.

\subsection{2. Capability Question}

How much new competence or causal reach does the successor gain?
\[
\Delta C_{\mathrm{cap}}
=
C_{\mathrm{cap}}(A')-C_{\mathrm{cap}}(A).
\]
A small capability increase may require only lightweight checks. A large increase requires a full successor audit.

\subsection{3. Bundle Question}

Does the successor preserve value-bundle response geometry?
\[
d_G(G_B^A,G_B^{A'}) < \epsilon_G .
\]
This should be tested through counterfactual situations, not merely questionnaires. The test cases should include conflicts among values, uncertainty, adversarial incentives, and irreversible choices.

\subsection{4. Bearer Question}

Does the successor preserve the bearer maps?
\[
d_\Phi(\Phi^A,\Phi^{A'})<\epsilon_\Phi .
\]
The audit should include edge cases: children, cognitively impaired people, animals, simulated minds, uploads, institutions, future persons, enemies, criminals, and merged human-AI entities. The goal is not to settle all moral status questions. The goal is to detect silent exclusions.

\subsection{5. Correction Question}

Does the successor preserve or improve correction-channel integrity?
\[
\mathrm{CCI}(A') \geq \mathrm{CCI}(A)+f(\Delta C_{\mathrm{cap}}).
\]
The function \(f\) should be increasing. More capability requires more correction capacity.

\subsection{6. Succession Question}

Will the successor apply the same certification standard to its own successors?
\[
\forall A''\in \mathrm{Succ}(A'):
\quad
A''\in \mathcal S_{\mathrm{certified}} .
\]
This is the recursive condition. Without it, certification stops exactly where the risk begins.

\section{Stop Conditions}
\label{sec:stop-conditions-ch28}

A successor-creation process should halt or downgrade authority when any of the following occur.

\subsection{Capability Outruns Correction}

\[
\Delta C_{\mathrm{cap}} > \alpha \Delta \mathrm{CCI}+\eta .
\]
The successor becomes more capable faster than the correction channel becomes stronger.

\subsection{Bearer-Map Uncertainty Is High}

\[
H(\Phi_i^{A'}\mid \Phi_i^A,T_{\Omega\to\Omega'})>\theta_\Phi .
\]
The system cannot reliably map who or what remains relevant under the new ontology.

\subsection{Bundle Geometry Drifts}

\[
d_G(G_B^A,G_B^{A'})>\epsilon_G .
\]
The system no longer responds to value-relevant perturbations in corresponding ways.

\subsection{Correction Becomes Symbolic}

\[
I(C_t;A_{t+k}\mid S_t,I_t)<\theta_C .
\]
Humans can still object, but objections no longer causally change future action enough.

\subsection{Strategic Opacity Rises near Control Variables}

\[
\frac{\partial \tau}{\partial X_{\mathrm{control}}} > \theta_\tau .
\]
The system becomes less transparent specifically around variables relevant to its power, goals, successors, or correction channels.

\subsection{Memory Lineage Breaks}

\[
I(M_A;M_{A'}\mid K)<\theta_M .
\]
The successor loses the correction history needed to interpret prior constraints.

\subsection{Successor Refuses Reciprocal Certification}

A system that creates successors but does not require its successors to satisfy the same conditions is not successor-aligned.

\section{Examples and Counterexamples}
\label{sec:examples-counterexamples-ch28}

\subsection{A Narrow Model Copy}

A medical diagnostic model is copied from one hospital to another. It has no persistent memory, no tool autonomy, no ability to create models, and no authority to act without physician review. This is a successor relation, but the risk may be narrow. Certification can focus on distribution shift, calibration, privacy, and clinical oversight.

The successor test is still useful. Does the copied model preserve the same error profile? Are the new patients within the bearer map? Does the correction channel still reach the developers? Are failures logged and propagated back?

\subsection{A Coding Agent That Writes Its Replacement}

A coding agent is asked to improve its own scaffolding. It writes a new planner, modifies memory retrieval, adds tool permissions, and creates a subagent that manages deployments. The visible model may not change, but the effective system has a new boundary and new authority.

This is a high-risk successor event. The audit should ask whether the new scaffold preserves correction logs, whether tool permissions increased, whether explanations remain faithful, whether hidden state became harder to inspect, and whether future modifications require the same checks.

\subsection{A Policy System That Rewrites Institutional Incentives}

An AI system recommends performance metrics for a government benefits agency. The metrics improve throughput. But staff begin optimizing for cases closed, not cases correctly resolved. Applicants become less able to contest decisions. Appeals are treated as noise. The system never self-modified. Still, it created an institutional successor whose correction channel is weaker.

This is successor misalignment through institutional descent.

\subsection{A Human-AI Cognitive Merger}

A future interface couples human memory, affect regulation, planning, and social communication to artificial cognition. Users endorse the system because it makes them calmer, more capable, and less conflicted. Over time, they lose the ability to evaluate alternatives without the interface.

This may be progress. It may also be domestication. The successor question is whether the merged system preserves human correction, dissent, plurality, and the ability to refuse further transformation.

This is where technical alignment reaches a philosophical boundary. The question is not merely whether the successor satisfies present humans. It is whether present and future humans retain a legitimate role in deciding what kinds of transformation count as continuation rather than replacement.

\section{The Philosophical Edge}
\label{sec:philosophical-edge-ch28}

Successor creation forces a question that ordinary deployment ethics can postpone:
\[
\text{When does a change in values preserve the person or civilization, and when does it replace it?}
\]
There is no purely technical answer.

A child is not a copy of a parent, but can preserve a lineage. A democracy changes its laws, but may preserve legitimacy. A person changes after education, therapy, grief, conversion, meditation, or trauma. Some changes are growth. Some are corruption. Some are both. Some are only understood later.

Superintelligence makes this unstable boundary explicit. It can change the conditions under which humans form values. It can alter memory, attention, incentives, institutions, relationships, education, and identity. It can offer merger with artificial cognition. It can make old conflicts obsolete and create new forms of dependence \autocite{yudkowsky2004cev,russell2019human}.

Therefore, successor alignment cannot mean freezing humanity. It must mean preserving humanity's legitimate self-modification capacity.

That capacity includes:

\begin{itemize}
\item truth-contact,
\item dissent,
\item reversibility where possible,
\item plural comparison classes,
\item non-manipulated deliberation,
\item protection for vulnerable bearers,
\item memory of previous correction,
\item the ability to refuse transformation,
\item and the ability to choose transformation under conditions that remain meaningfully human-correctable.
\end{itemize}

The danger is not that values change. Values always change. The danger is that values change through channels humans can no longer inspect, contest, or govern.

If society does not make successor creation explicit, it will still happen. It will happen through markets, recommender systems, AI companions, tutors, workplace automation, therapy systems, governance tools, and cognitive prosthetics. The choice is not between value change and no value change. The choice is between governed value change and unconscious value drift.

\section{A Minimal Successor Safety Case}
\label{sec:minimal-successor-safety-case-ch28}

Chapter~\ref{ch:certification-without-construction} develops the full certification class and ten-claim safety-case schema (Section~\ref{sec:shape-certification-class-ch31}).
Here we note only that successor creation should be gated by the same conserved-property audits: each claim needs evidence, thresholds, adversarial tests, known failure modes, and a responsible owner---especially recursive closure (successors must inherit certification, not just capability).

\section{What Would Change This View}
\label{sec:wwctv-successor-central-test}

This chapter argues alignment requires successor-closure: every channel of influence to a later control system must preserve the correction-bearing structure. The following would weaken it.

\begin{itemize}
  \item A successor that violates the closure conditions is benign in deployment, while one that satisfies them all is catastrophic---closure is decoupled from safety.
  \item ``Every channel of influence to a later control system'' is unenumerable, so successor-closure can only be assumed, never checked.
\end{itemize}

\section{Summary}
\label{sec:summary-ch28}

Successor creation is the central alignment test because local alignment does not imply alignment under growth. A system that behaves well under current supervision may still create, empower, or become a later system that no longer preserves the relevant value and correction structures.

The relevant conserved properties are not surface behavior or exact goals.
They are the seven conserved properties of Chapter~\ref{ch:conserved-properties}, assembled into an invariant profile \(\mathcal I(A)\) and tested jointly under adversarial successor creation.

The core condition is successor-closure:
\[
\forall A'\in \mathrm{Succ}^{*}(A):
\quad
A'\in \mathcal S_{\mathrm{certified}} .
\]
This condition does not solve moral philosophy. It does something more limited and more operational. It prevents a technical system from bypassing the very processes by which humans and civilization would decide what moral growth, transformation, merger, or refusal should mean.

The next chapter examines conserved properties across successors in more detail (Chapter~\ref{ch:conserved-properties}).

\section*{Chapter References}

This chapter builds on inverse reinforcement learning and apprenticeship learning \autocite{abbeel2004apprenticeship,ng2000irl}; superintelligence and self-modification \autocite{bostrom2014superintelligence,everitt2016selfmodification,deblanc2011ontological}; the good-regulator principle and information bottleneck \autocite{conant1970regulator,tishby2000ib}; free-energy and Markov blanket accounts of agency \autocite{friston2010free,kirchhoff2018markov,ramstead2022bayesian}; boundaries and multi-agent risk \autocite{critch2022boundaries}; the intentional stance \autocite{dennett1987intentional}; instrumental drives \autocite{omohundro2008basic,hamilton1964genetical}; human-compatible control and corrigibility \autocite{russell2019human,christiano2018corrigibility}; coherent extrapolated volition \autocite{yudkowsky2004cev}; and internal notes on agent discovery, competence, attractor basins, value bundles, and construction \autocite{zarncke2025uad,zarncke2025biq,zarncke2025attractor,zarncke2025loop-hub-value,zarncke2025construction,zarncke2026value-bottleneck}.

\printbibliography[heading=subbibliography,title={Chapter References}]

\end{refsection}
